\section{Related Work}

\subsection{Direct Preference Optimization}

Direct Preference Optimization (DPO) provides a closed-form solution for learning from human preferences without training a separate reward model \cite{rafailov2023direct}. Given paired preferences $(y_w, y_l)$ for prompt $x$, DPO optimizes:
\begin{equation}
\mathcal{L}_{\text{DPO}} = -\log \sigma\left(\beta \log \frac{\pi_\theta(y_w|x)}{\pi_{\text{ref}}(y_w|x)} - \beta \log \frac{\pi_\theta(y_l|x)}{\pi_{\text{ref}}(y_l|x)}\right)
\end{equation}

This formulation has become the dominant approach for preference-based fine-tuning due to its simplicity and stability. However, DPO optimizes a single objective: matching human preferences. Our work asks whether auxiliary objectives can capture dimensions of response quality orthogonal to preference labels.

\subsection{Multi-Objective Extensions to DPO}

Several works extend DPO with additional objectives. MODPO introduces margin-based auxiliary losses that provide additional training signal beyond binary preferences \cite{zhou2024modpo}. PAMA (Pareto Multi-Objective Alignment) frames alignment as multi-objective optimization across multiple preference dimensions \cite{he2025pama}. These methods share a common pattern: they add objectives derived from preference data. BiDPO differs by adding an objective orthogonal to preference labels.

\subsection{Human Agency in AI Systems}

Mitelut et al.\ argue that intent-aligned AI may inadvertently deplete human agency by optimizing for immediate task completion rather than capability transfer \cite{mitelut2023intent}. This theoretical framework motivates our empirical investigation: can we design training objectives that preserve agency-related patterns?
